\documentclass{article}
\usepackage[UTF8]{ctex}
\usepackage{amsmath,amssymb,amsfonts}
\usepackage{enumitem}
\usepackage{geometry}
\geometry{a4paper, margin=1in}
% === 错题本排版优化（不参与题目标识） ===
\setlength{\parindent}{0pt}
\setlength{\parskip}{0.35em plus 0.10em minus 0.08em}
\setlength{\emergencystretch}{3em}
\setlist{itemsep=0.25em, topsep=0.40em, parsep=0pt}
\allowdisplaybreaks[2]
% === 错题本排版优化结束 ===
\begin{document}
\section*{北京邮电大学2026年高等代数真题}
\section*{13．北京邮电大学2026年高等代数与概率论试卷}

\subsection{题目1}
一．设 \(A, B\) 是 \(n\) 阶实方阵，\(A\) 可对角化，且 \(B A^{2}=A+B\) 。证明：

\begin{itemize}
  \item[（1）] 1 不是 \(A\) 的特征值．
  \item[（2）] 存在可逆矩阵 \(P\) ，使得 \(P^{-1} A P, P^{-1} B P\) 均为对角矩阵。
\end{itemize}

\subsection{题目2}
二．求一个矩阵 \(A\) ，使得它的伴随矩阵是 \(A^{*}=\left(\begin{array}{lll}1 & 1 & 1 \\ 2 & 2 & 2 \\ 3 & 3 & 3\end{array}\right)\) ．

\subsection{题目3}
三。设整系数多项式 \(f(x)\) 的次数是 \(n=2 m\) 或 \(n=2 m+1\left(m \in \mathbb{N}^{+}\right)\)，证明：如果有 \(k(\geq 2 m+1)\) 个不同的整数 \(a_{1}, a_{2}, \cdots, a_{k}\) ，使得每个 \(f\left(a_{l}\right)\) 取值为 1 或 -1 ，那么 \(f(x)\) 在有理数域上不可约．

\subsection{题目4}
四．设 \(A, B\) 分别是数域 \(P\) 上的 \(m \times n, n \times s\) 矩阵，\(V=\left\{B \gamma \mid \gamma \in P^{n}, A B \gamma=0\right\}\) 是 \(n\) 维向量空间\(P^{n}\) 的子空间，证明：维 \((V)=\) 秩 \((B)-\) 秩 \((A B)\) 。

\subsection{题目5}
五。设 \(V\) 是 \(n\) 维欧氏空间，\(\sigma\) 是 \(V\) 上的正交变换，设子空间

\[
V_{1}=\{\alpha \mid \sigma(\alpha)=\alpha, \alpha \in V\}, V_{2}=\{\beta \mid \beta=\alpha-\sigma(\alpha), \alpha \in V\} .
\]

证明：\(V=V_{1} \oplus V_{2}\) ．

\subsection{题目6}
六。独立重复进行伯努利实验，假设每次实验结果为 \(A\) 和 \(\bar{A}\) ，且满足 \(P(A)=p \in(0,1)\) ，设 \(x_{k}\) 表示实验结果 \(A\) 第 \(k\) 次出现所需的试验次数，\(k=1,2, \cdots\) ，令 \(m, n\) 为正整数。

\begin{itemize}
  \item[（1）] 求概率 \(P\left(x_{1}=n, x_{3}=m+n+1\right)\) 。
  \item[（2）] 求条件概率 \(P\left(x_{1}=m \mid x_{2}=m+n\right)\) ．
  \item[（3）] 求数学期望 \(E\left[x_{k}\right]\) 和方差 \(D\left[x_{k}\right], k=1,2, \cdots\) ．
\end{itemize}

\subsection{题目7}
七．设二维随机变量 \((X, Y)\) 的概率密度函数为

\[
f(x, y)= \begin{cases}c(x+y), & 0 \leq y \leq x \leq 1, \\ 0, & \text { 其他. }\end{cases}
\]

\(c\) 为待定系数．

\begin{itemize}
  \item[（1）] 求 \(c\) 。
  \item[（2）] 求 \(X, Y\) 的边缘密度函数．
  \item[（3）] 求协方差 \(\operatorname{cov}(X, Y)\) ．
\end{itemize}

\subsection{题目8}
八．设 \(f(x, y)=\left\{\begin{array}{ll}y e^{-y(1+x)}, & x>0, y>0, \\ 0, & \text { 其他．}\end{array}\right.\) 令 \(\left\{\begin{array}{l}U=X Y, \\ V=Y .\end{array}\right.\)

\begin{itemize}
  \item[（1）] 求 \((U, V)\) 的概率密度函数 \(g(u, v)\) ．
  \item[（2）] 求 \(U, V\) 的边缘概率密度函数 \(g_{U}(u)\) 和 \(g_{V}(v)\) ，并判别 \(U, V\) 的独立性．
  \item[（3）] 求 \(Z=U+V\) 的概率密度函数．
\end{itemize}

\subsection{题目9}
九．设随机变量 \(X\) 的分布律为 \(P\{X=k\}=\frac{1}{2^{k}}, k=1,2, \cdots\) ，令 \(Y=\sin \left(\frac{\pi}{2} X\right)\) ．

\begin{itemize}
  \item[（1）] 求 \(Y\) 的分部律．
  \item[（2）] 设随机变量序列 \(\left\{Y_{n}, n=1,2, \cdots\right\}\) 独立同分布，且与 \(Y\) 有相同的分布函数．
  \item[] （2－1）对于任意的 \(\varepsilon>0\) ，利用切比雪夫不等式证明：
\[
\lim _{n \rightarrow \infty} P\left(\left|\frac{1}{n} \sum_{k=1}^{n} Y_{k}-\frac{2}{5}\right|>\varepsilon\right)=0
\]
  \item[] （2－2）设常数 \(a>0\) ，满足 \(\Phi\left(\sqrt{\frac{75}{38}}\right)=a\) ，其中 \(\Phi\) 为标准正态分布的分布函数，求
\[
\lim _{n \rightarrow \infty} P\left(\frac{1}{\sqrt{n}}\left|\sum_{k=1}^{n} Y_{k}-\frac{2 n}{5}\right|>1\right)
\]
  \item[十．] 设随机变换 \(X, Y\) 相互独立，概率密度函数分别为
\[
f_{X}(x)=\left\{\begin{array}{ll}
\lambda e^{-\lambda x}, & x>0, \\
0, & \text { 其他. }
\end{array} \quad f_{Y}(y)=\left\{\begin{array}{ll}
\mu e^{-\mu y}, & y>0, \\
0, & \text { 其他. }
\end{array} \quad \lambda, \mu>0 .\right.\right.
\]
  \item[（1）] 求 \(P(X>Y)\) ．
  \item[（2）] 令 \(M=\max \{X, Y\}, N=\min \{X, Y\}\) ，求其概率密度函数 \(f_{M}(x), f_{N}(x)\) ．
  \item[（3）] 求 \(N\) 的特征函数 \(\varphi_{N}(t)\) ．
\end{itemize}
\end{document}
